\chapter{Anexos de trazabilidad y reproducibilidad}
\begingroup
\sloppy

\section{Matriz objetivo--evidencia}

La Tabla~\ref{tab:matriz-objetivos-anexo} resume el estado verificable al 2 de
septiembre de 2026. ``Parcial'' no significa incumplimiento oculto: identifica
el producto disponible y la evidencia que todavía falta para cerrar el
objetivo en el prototipo físico.

\begin{longtable}{p{0.09\textwidth}p{0.27\textwidth}p{0.31\textwidth}p{0.22\textwidth}}
\caption{Matriz resumida de objetivos, evidencia y estado.}\label{tab:matriz-objetivos-anexo}\\
\toprule
Objetivo & Resultado demostrable & Evidencia principal & Estado y faltante \\
\midrule
\endfirsthead
\toprule
Objetivo & Resultado demostrable & Evidencia principal & Estado y faltante \\
\midrule
\endhead
OE1 & Modelo cinemático y dinámico nominal implementado, consistente y sometido a sensibilidad. & Capítulo 7, pruebas de modelo y gráficas cinemáticas. & Parcial avanzado; faltan geometría, masas, límites y contraste físico. \\
OE2 & Arquitectura ROS~2, sensores simulados y supervisor con nueve escenarios provocados aprobados. & JSON de seguridad, código, Raspberry y PCA9685. & Parcial avanzado; faltan OE, protección, rearme y validación física. \\
OE3 & Gateo y marcha paso repetibles en simulación; precarga 2,0 repetida con vuelo trasero sostenido. & Bolsas, análisis por ciclo y validación Gazebo--MuJoCo. & Parcial avanzado; falta marcha nominal física. \\
OE4 & Cinco semillas PPO entrenadas y evaluadas; ninguna escala positiva fue aceptada. & Curvas, modelos, bolsas y diagnóstico de escala cero. & Ejecutado sin mejora aceptada; requiere rediseño antes de repetir. \\
OE5 & Comparación descriptiva corrigió una referencia inválida y bloqueó transferencia. & CSV y figuras de comparación nominal--PPO. & Parcial; falta campaña emparejada final de cinco ensayos por condición y marcha. \\
\bottomrule
\end{longtable}

La matriz viva, con criterio de cierre y rutas completas, se conserva en
\path{Documento_TESIS/MATRIZ_OBJETIVO_EVIDENCIA.md}.

\section{Protocolos aplicables}

Los procedimientos no se transcriben de forma abreviada para evitar dos
versiones divergentes. Las fuentes versionadas son:

\begin{itemize}
\item \path{Documentacion/PROTOCOLO_EXPERIMENTAL_BORRADOR.md}: ensayo, ciclo,
fallo, intervención, métricas y criterios provisionales;
\item \path{Documentacion/PROTOCOLO_CARACTERIZACION_FISICA.md}: fotografías,
geometría, masas y mapa de canales;
\item \path{Documentacion/PROTOCOLOS_FISICOS/}: puesta en marcha, calibración,
identificación y progresión de la marcha;
\item \path{Experimentos/PROTOCOLO_NODO_PRUEBAS_SEGURIDAD.md}: estímulos
aislados del supervisor;
\item \path{Experimentos/PROTOCOLO_PERTURBACIONES_GAZEBO.md}: condiciones para
fricción, empujes, ruido y retardos;
\item \path{Experimentos/PROTOCOLO_PPO_RESIDUAL.md}: contrato y campaña de la
capa aprendida.
\end{itemize}

\section{Comandos mínimos de reproducción}

La construcción y las pruebas deterministas se ejecutan desde la raíz:

\begin{verbatim}
source /opt/ros/humble/setup.bash
colcon build --packages-select \
  nova_sm3_description nova_gait_controller
source install/setup.bash
python3 -m pytest -q src/nova_gait_controller/test
\end{verbatim}

Los nueve escenarios dinámicos del supervisor se reproducen con:

\begin{verbatim}
ROS_DOMAIN_ID=93 \
Experimentos/ejecutar_pruebas_dinamicas_supervisor.sh \
  /tmp/pruebas_dinamicas_supervisor
\end{verbatim}

La línea base se lanza con \texttt{demo.launch.py}; antes de grabar se debe
comprobar instancia única mediante \texttt{ros2 node list | sort | uniq -c} y
confirmar los suscriptores del grabador en el tópico de órdenes. El análisis se
realiza con \path{Experimentos/analizar_gateo_rosbag.py} y
\path{Experimentos/analizar_contactos_rosbag.py}. La equivalencia se agrega con
\path{Experimentos/comparar_simuladores_equivalentes.py}; en MuJoCo se reinicia
primero el keyframe \texttt{stand}. Las rutas exactas de bolsa,
configuración y hashes se toman del registro de cada ensayo, nunca de una bolsa
con marca \texttt{ENSAYO\_INVALIDO.md}.

\section{Reglas de integridad y repetición}

Cada resultado debe conservar fecha, versión, parámetros, marcadores de inicio
y fin, número de ciclos, criterio de exclusión y SHA-256. El ciclo 1 se mantiene
como transitorio; la unidad independiente es el ensayo. No se mezclan bolsas
creadas con cadencias, geometrías o perfiles diferentes. Las tentativas fallidas
se conservan con su causa y no entran en las estadísticas aceptadas.

El flujo de integración continua en
\path{.github/workflows/pruebas-automaticas.yml} compila los dos paquetes,
ejecuta las pruebas Python y repite los escenarios dinámicos en cada
\texttt{push} y \texttt{pull\_request}. Esto asegura regresión de software, no
reemplaza campañas dinámicas extensas ni validación física.

\section{Lista de transición al prototipo}

El orden F0--F4 se conserva en el plan de validación física fechado el 2 de
septiembre y versionado en la carpeta \path{Raspberry}. La
secuencia obligatoria es: ensamble y caracterización sin energía; seguridad
eléctrica y OE; calibración canal por canal; integración suspendida y luego con
soporte; instrumentación y contraste del modelo. Ninguna etapa autoriza PPO o
marcha física si la anterior permanece abierta.
\endgroup
